\documentclass[11pt,a4paper]{report}

\usepackage[margin=22mm]{geometry}

\usepackage{amsmath,amssymb,bm}

\usepackage{graphicx}

\usepackage{booktabs,longtable,array}

\usepackage[hidelinks]{hyperref}

\usepackage{float}

\usepackage{caption}

\setlength{\parskip}{0.5em}

\setlength{\parindent}{0pt}

\title{Inverse Design and Sequential Dual-SLM Realisation of a Source-Scale Hexagonal Vector Bessel Beam}

\author{Nathan source-scale branch technical evidence draft\\Generated by Codex on 2026-07-09}

\date{}

\begin{document}

\maketitle

\begin{center}\textbf{Scope statement:} This report covers the source-scale branch. It does not claim experimental replication or microfabrication/sample-plane success.\end{center}

\tableofcontents

\listoffigures

\chapter*{Abstract}

\addcontentsline{toc}{chapter}{Abstract}

This technical evidence draft derives the Nathan source-scale hexagonal vector Bessel beam model, traces the ideal and realistic synthesis routes, documents inverse reconstruction, audits sampling and strict-hexagon classification, and specifies the corrected sequential single-beam dual-SLM laboratory architecture. The current conclusion is that the source-scale sequential architecture is validated in simulation and sufficiently specified to justify a source-scale laboratory trial. Experimental replication and microfabrication/sample-plane success remain separate unresolved stages.

\chapter*{Notation and Abbreviations}

\addcontentsline{toc}{chapter}{Notation and Abbreviations}

\begin{longtable}{p{0.18\textwidth}p{0.72\textwidth}}

\toprule

Symbol or term & Meaning\\

\midrule

\endfirsthead

\toprule

Symbol or term & Meaning\\

\midrule

\endhead

\(A(r)\) & Gaussian source amplitude envelope.\\

\(w_0\) & Source-scale beam radius used by the Nathan parity branch.\\

\(\theta\) & Transverse-plane angular coordinate, \(\operatorname{atan2}(y,x)\).\\

\(\alpha(\theta)\) & Local linear-polarisation angle after the radial/azimuthal sector offset.\\

SLM & Spatial light modulator; here a phase-only LCOS panel.\\

HWP & Half-wave plate; a retarder used for polarisation rotation or channel swapping.\\

QWP & Quarter-wave plate; the final retarder that maps the dual phase channels into the target linear vector field.\\

PBS & Polarising beamsplitter. It appears only as superseded split-arm context here, not in the accepted canonical architecture.\\

ASM & Angular-spectrum method for free-space propagation.\\

4F & Fourier filtering relay with two lenses separated by focal lengths, used to select the carrier-shifted +1 diffraction order.\\

\(S_0,S_1,S_2,S_3\) & Stokes observables used to inspect intensity and polarisation state.\\

\bottomrule

\end{longtable}

\chapter{Introduction}

The project asks whether Nathan's segmented radial/azimuthal source-scale vector Bessel field can be reproduced with a programmable laboratory route. The answer at this stage is yes in simulation for the source-scale branch. The aim is not to claim microfabrication or sample-plane success. It is to establish the mathematical field, the optical transformations, the realistic carrier/4F implementation, and the corrected sequential dual-SLM bench architecture.

\textbf{Key rule:} the final architecture is sequential and collinear. No canonical PBS split/recombine interferometer arms are used.

\chapter{Physical and Mathematical Background}

The source field begins as a Gaussian amplitude A(r)=exp(-r\textasciicircum{}2/w0\textasciicircum{}2). The six-sector structure is not a scalar intensity mask. It lives in the local polarization angle alpha(theta)=theta+phi0(theta), where phi0 alternates between radial and azimuthal sectors. Thus the pre-axicon target is E=A[cos alpha, sin alpha]\textasciicircum{}T, and S0 remains Gaussian before the axicon. Stokes fields reveal the polarization structure; intensity alone does not.

\begin{align}A(r)&=\exp(-r^2/w_0^2)\\ \alpha(\theta)&=\theta+\phi_0(\theta)\\ \mathbf E_\mathrm{target}&=A(r)[\cos\alpha,\sin\alpha]^T\end{align}

\begin{figure}[H]

\centering

\includegraphics[width=0.92\textwidth]{\detokenize{evidence_pack/equations/F04_m2p_target_alpha_sector_highres.png}}

\caption{M2P target alpha and sector map.}

\label{fig:F04}

\end{figure}

\begin{figure}[H]

\centering

\includegraphics[width=0.92\textwidth]{\detokenize{evidence_pack/equations/F05_m2p_stokes_comparison_highres.png}}

\caption{M2P pre-axicon Stokes comparison.}

\label{fig:F05}

\end{figure}

\chapter{V0 Source-Model Reproduction}

V0 established the reference observable and grid convention. The Nathan-style axis-sampled grid is required because a zero-straddling grid creates a central artefact at the angular singularity. The literal source port and the project implementation reproduce the same source-scale field. This is the reference used by later route comparisons.

\begin{figure}[H]

\centering

\includegraphics[width=0.92\textwidth]{\detokenize{evidence_pack/source_validation/F01_nathan_visual_ladder_v0_reference_vs_reproduction.png}}

\caption{V0 source parity and reference/reproduction comparison.}

\label{fig:F01}

\end{figure}

\begin{figure}[H]

\centering

\includegraphics[width=0.92\textwidth]{\detokenize{evidence_pack/source_validation/F02_nathan_source_observable_audit_grid_convention.png}}

\caption{Grid convention audit showing why Nathan axis-sampled centring is required.}

\label{fig:F02}

\end{figure}

\begin{figure}[H]

\centering

\includegraphics[width=0.92\textwidth]{\detokenize{evidence_pack/source_validation/F03_v0_reference_crop_highres.png}}

\caption{High-resolution V0 reference crop and propagated structure.}

\label{fig:F03}

\end{figure}

\chapter{Vector Axicon Propagation}

After the target vector field reaches the axicon, the local radial and azimuthal components see the source-scale conical phase and p/s response. The propagated observable is total intensity |Ex|\textasciicircum{}2+|Ey|\textasciicircum{}2+|Ez|\textasciicircum{}2. The vector angular-spectrum propagator enforces transversality and advances each spectral component by exp(i kz z). The hexagonal structure is therefore a downstream propagation result, not merely a pre-axicon drawing.

\chapter{Ideal Jones-Chain Synthesis}

M2P showed two analytic routes. A patterned HWP with beta=alpha/2 maps horizontal input to the target Jones vector. The dual-SLM/QWP route uses phase masks phi\_H=+alpha and phi\_V=-alpha+pi/2, followed by a QWP in the selected project convention. The derivation is power preserving; relative piston rotates the local polarization state and remains visible to polarimetry even when total intensity is insensitive.

\begin{figure}[H]

\centering

\includegraphics[width=0.92\textwidth]{\detokenize{evidence_pack/equations/F04_m2p_target_alpha_sector_highres.png}}

\caption{M2P target alpha and sector map.}

\label{fig:F04}

\end{figure}

\begin{figure}[H]

\centering

\includegraphics[width=0.92\textwidth]{\detokenize{evidence_pack/equations/F05_m2p_stokes_comparison_highres.png}}

\caption{M2P pre-axicon Stokes comparison.}

\label{fig:F05}

\end{figure}

\chapter{Forward Source-Scale Bench Replication}

M2N propagated the ideal patterned-HWP, ideal dual-SLM/QWP, and realistic dual-SLM + carrier + 4F + QWP routes through the same source axicon. The ideal routes reproduce V0 at z=60 mm. The realistic carrier/4F route keeps the strict visual hexagon with z=60 correlation around 0.9936 and first-order efficiency around 0.949.

\begin{figure}[H]

\centering

\includegraphics[width=0.92\textwidth]{\detokenize{evidence_pack/sequential_architecture/F07_fig3A_v0_ideal_realistic_sequential.png}}

\caption{M2W-FIX V0 versus ideal sequential versus realistic sequential comparison.}

\label{fig:F07}

\end{figure}

\chapter{Backward / Adjoint Mask Synthesis}

M2Q starts from the full complex V0 target at z=60 mm, backpropagates it, applies an inverse axicon operation, and removes the QWP. The recovered H/V phase solution independently returns phi\_H=+alpha and phi\_V=-alpha+pi/2. This is important because it confirms the forward analytic mask convention from the inverse problem.

\begin{figure}[H]

\centering

\includegraphics[width=0.92\textwidth]{\detokenize{evidence_pack/inverse_design/F06_m2q_phase_only_masks_highres.png}}

\caption{M2Q recovered phase-only masks and inverse design diagnostics.}

\label{fig:F06}

\end{figure}

\chapter{Realistic 4F Filtering}

A displayed phase mask includes the desired phase plus a carrier 2*pi*nu*x. In the Fourier plane the +1 order is displaced by x=lambda f nu. With lambda=1029 nm, f=300 mm and nu=6.25 lp/mm, the displacement is about 1.929 mm and the recommended iris diameter is about 1.54 mm. The common 4F route filters both polarization channels together.

\begin{figure}[H]

\centering

\includegraphics[width=0.92\textwidth]{\detokenize{evidence_pack/sequential_architecture/F07_fig3A_v0_ideal_realistic_sequential.png}}

\caption{M2W-FIX V0 versus ideal sequential versus realistic sequential comparison.}

\label{fig:F07}

\end{figure}

\begin{figure}[H]

\centering

\includegraphics[width=0.92\textwidth]{\detokenize{evidence_pack/hardware/F09_fig2_target_and_sequential_masks.png}}

\caption{Native SLM1/SLM2 masks and target vector-field diagnostics.}

\label{fig:F09}

\end{figure}

\chapter{Lab Realism and Tolerance}

M2S introduced quantisation, fill factor, H/V imbalance, piston, QWP errors, iris errors, registration, axicon decentre and z offsets. Within tested ranges, most representative imperfections were relatively forgiving. The dominant practical sensitivity is hologram/mask centre relative to the axicon axis. A 0.5 mm axicon/mask offset is not acceptable blindly but can be substantially recovered by digital recentring.

\begin{figure}[H]

\centering

\includegraphics[width=0.92\textwidth]{\detokenize{evidence_pack/realism/F10_fig3B_clean_moderate_bad_realism.png}}

\caption{Realism degradation: clean, moderate and bad cases.}

\label{fig:F10}

\end{figure}

\begin{figure}[H]

\centering

\includegraphics[width=0.92\textwidth]{\detokenize{evidence_pack/correction/F11_fig3C_degraded_corrected_recovery.png}}

\caption{Corrected recovery from a 0.5 mm axicon/mask offset.}

\label{fig:F11}

\end{figure}

\chapter{High-Resolution and Sampling Audit}

M2U separated numerical resolution from display resolution. N=512 is Nyquist-safe but not publication-preferred. N=1024 is useful confirmation. N=1536 is publication-selected for the hero comparison with about 9.89 native samples per radial fringe. M2W-FIX adds SAS-scaled focus crops so close-up panels are rendered on finer output grids while classifier metrics remain tied to native validated arrays.

\begin{figure}[H]

\centering

\includegraphics[width=0.92\textwidth]{\detokenize{evidence_pack/source_validation/F03_v0_reference_crop_highres.png}}

\caption{High-resolution V0 reference crop and propagated structure.}

\label{fig:F03}

\end{figure}

\begin{figure}[H]

\centering

\includegraphics[width=0.92\textwidth]{\detokenize{evidence_pack/sequential_architecture/F07_fig3A_v0_ideal_realistic_sequential.png}}

\caption{M2W-FIX V0 versus ideal sequential versus realistic sequential comparison.}

\label{fig:F07}

\end{figure}

\begin{figure}[H]

\centering

\includegraphics[width=0.92\textwidth]{\detokenize{evidence_pack/source_validation/F12_fig4A_transverse_evolution.png}}

\caption{Transverse propagation with physical context and SAS-scaled focus crops.}

\label{fig:F12}

\end{figure}

\chapter{Optimisation Integrity and Strict Hexagon Gate}

M2U2-FIX repaired an important failure: high full-field correlation can be dominated by dark background and can hide fourfold or triangular structures. The old best compromise is permanently forbidden. The strict gate uses class, focus support, angular correlation, c60/c120 behaviour, dark-core limits and fourfold vetoes. Its 0.997 reference-correlation floor is a project-calibrated eligibility threshold, not a universal definition of a hexagon.

\begin{figure}[H]

\centering

\includegraphics[width=0.92\textwidth]{\detokenize{evidence_pack/provenance/F17_hexagon_classifier_truth_table_highres.png}}

\caption{Strict classifier truth table and gate calibration.}

\label{fig:F17}

\end{figure}

\begin{figure}[H]

\centering

\includegraphics[width=0.92\textwidth]{\detokenize{evidence_pack/superseded/S01_best_compromise_m2u2_opt_003_c5.75_i0.32_q-0.25_r+0.0_p0.10_audit.png}}

\caption{Forbidden old best-compromise audit figure.}

\label{fig:S01}

\end{figure}

\chapter{Native-Panel and Hardware Provenance}

The panel raster is 1920 x 1080 at 8 um pitch, giving a 15.36 x 8.64 mm active area. The exact PLUTO-2.1 NIR-149 identity is externally supplied until physical labels/manuals are read. Phase stroke and LUT are calibration-required. Camera scale, z-stage and QWP mount sign also remain bench calibrations.

\begin{figure}[H]

\centering

\includegraphics[width=0.92\textwidth]{\detokenize{evidence_pack/hardware/F09_fig2_target_and_sequential_masks.png}}

\caption{Native SLM1/SLM2 masks and target vector-field diagnostics.}

\label{fig:F09}

\end{figure}

\chapter{Sequential Single-Beam Dual-SLM Architecture}

M2W-FIX supersedes the original split-arm diagram. The accepted implementation is collinear and sequential: equal H/V preparation, SLM1, optional swap HWP, SLM2, optional swap-back, common 4F, QWP, axicon and camera. The sequential Jones derivation gives E4=A/sqrt(2)[exp(i phi\_H), exp(i phi\_V)]\textasciicircum{}T. The pre-QWP and post-QWP overlaps are 1.0, the ideal z=60 correlation is 1.0, and the realistic sequential route remains strict-eligible.

\begin{align}E_0&=\frac{A}{\sqrt2}[1,1]^T\\E_1&=\frac{A}{\sqrt2}[e^{i\phi_H},1]^T\\E_2&=\frac{A}{\sqrt2}[1,e^{i\phi_H}]^T\\E_3&=\frac{A}{\sqrt2}[e^{i\phi_V},e^{i\phi_H}]^T\\E_4&=\frac{A}{\sqrt2}[e^{i\phi_H},e^{i\phi_V}]^T.\end{align}

The masks are \(\phi_H=+\alpha+\) carrier and \(\phi_V=-\alpha+\pi/2+\) carrier.

\begin{figure}[H]

\centering

\includegraphics[width=0.92\textwidth]{\detokenize{evidence_pack/sequential_architecture/F08_fig1_sequential_optical_architecture.png}}

\caption{Corrected sequential single-beam optical system.}

\label{fig:F08}

\end{figure}

\begin{figure}[H]

\centering

\includegraphics[width=0.92\textwidth]{\detokenize{evidence_pack/sequential_architecture/F07_fig3A_v0_ideal_realistic_sequential.png}}

\caption{M2W-FIX V0 versus ideal sequential versus realistic sequential comparison.}

\label{fig:F07}

\end{figure}

\chapter{Sequential Power / Energy Ledger}

The corrected power model has no split-arm H/V or PBS recombination rows. It tracks one beam through preparation, SLM1, optional swap, SLM2, common 4F selection, QWP, axicon, z=60 integration and useful-region power. The 1 W and 10 W examples are linear model examples only and are not damage-threshold claims.

\begin{figure}[H]

\centering

\includegraphics[width=0.92\textwidth]{\detokenize{evidence_pack/hardware/F14_fig4C_sequential_power_flow.png}}

\caption{Sequential-route power flow.}

\label{fig:F14}

\end{figure}

\chapter{Closed-Loop Correction}

The camera observes final intensity, centre, symmetry, dark core, lobe balance and z structure. Shack-Hartmann measurements address low-order common wavefront errors. Polarimetry checks pre-axicon vector field, H/V balance, relative phase and QWP sign. The bounded correction loop improves several degraded cases, but not every failure is fully recovered to strict eligibility.

\begin{figure}[H]

\centering

\includegraphics[width=0.92\textwidth]{\detokenize{evidence_pack/correction/F11_fig3C_degraded_corrected_recovery.png}}

\caption{Corrected recovery from a 0.5 mm axicon/mask offset.}

\label{fig:F11}

\end{figure}

\begin{figure}[H]

\centering

\includegraphics[width=0.92\textwidth]{\detokenize{evidence_pack/correction/F16_fig5B_combined_correction_results.png}}

\caption{Combined realism and correction results.}

\label{fig:F16}

\end{figure}

\chapter{Final Lab Build}

The recommended bench is a 1029 nm Gaussian source, POL/HWP preparation, sequential SLM1/SLM2 masks, common 300 mm 4F with 6.25 lp/mm carrier and 1.54 mm iris, QWP at nominal code -45 deg, 2 deg n=1.458 axicon, and a camera z scan around z=60 mm. The six-piece segmented optic is not required for this programmable source-scale trial.

\begin{figure}[H]

\centering

\includegraphics[width=0.92\textwidth]{\detokenize{evidence_pack/sequential_architecture/F08_fig1_sequential_optical_architecture.png}}

\caption{Corrected sequential single-beam optical system.}

\label{fig:F08}

\end{figure}

\begin{figure}[H]

\centering

\includegraphics[width=0.92\textwidth]{\detokenize{evidence_pack/source_validation/F13_fig4B_xz_yz_z_diagnostics.png}}

\caption{x-z and y-z propagation maps plus z diagnostics.}

\label{fig:F13}

\end{figure}

\chapter{Conclusions}

Nathan's source-scale field has been reproduced in simulation; the propagated result is a genuine strict-gate hexagonal field; the vector field can be synthesised analytically; inverse propagation recovers the same masks; realistic common-4F filtering preserves the beam; the main practical sensitivity is mask/axicon centring; the sequential single-beam dual-SLM architecture is valid in simulation. Source-scale lab trial is justified. Microfabrication/sample-plane success is not established.

\chapter{Limitations}

The exact SLM phase stroke and per-panel LUT are not measured here. Camera scale and z-stage calibration remain unresolved. Exact LC-director orientation is unverified. The 4F geometry is nominal until bench measurement. Simulations do not prove experimental success. Strict thresholds are project-calibrated. No microfabrication/sample-plane success claim is made.

\chapter{Next Steps}

Read physical SLM labels and determine LC directors; measure phase LUTs; deploy native masks; calibrate 4F displacement and iris; validate pre-axicon Stokes fields; insert the axicon; perform camera z scans; digitally recentre; add Shack-Hartmann correction; then use a bounded camera-driven correction loop before revisiting microfabrication adaptation.

\appendix

\chapter{Equation Registry}

\begin{longtable}{p{0.18\textwidth}p{0.18\textwidth}p{0.18\textwidth}p{0.18\textwidth}p{0.18\textwidth}}
\caption{Machine-readable equation registry excerpt.}\label{tab:eq-registry}\\
\toprule
equation\_id & name & latex & code & validation\_status\\
\midrule
\endfirsthead
\toprule
equation\_id & name & latex & code & validation\_status\\
\midrule
\endhead
EQ01 & Gaussian source amplitude & A(r)=\textbackslash{}exp(-r\textasciicircum{}2/w\_0\textasciicircum{}2) & vbb\_study/digital\_twin/nathan\_vector\_hexagon.py::mode2n\_source\_target & validated\\
EQ02 & Angular coordinate & \textbackslash{}theta=\textbackslash{}operatorname\{atan2\}(y,x) & nathan\_alpha\_map / source\_parity\_grid & validated\\
EQ03 & Six-sector orientation angle & \textbackslash{}alpha(\textbackslash{}theta)=\textbackslash{}theta+\textbackslash{}phi\_0(\textbackslash{}theta),\textbackslash{}quad \textbackslash{}phi\_0\textbackslash{}in\textbackslash{}\{0,\textbackslash{}pi/2\textbackslash{}\} & nathan\_alpha\_map & validated\\
EQ04 & Target Jones vector & \textbackslash{}mathbf E\_\textbackslash{}mathrm\{target\}=A(r)[\textbackslash{}cos\textbackslash{}alpha,\textbackslash{} \textbackslash{}sin\textbackslash{}alpha]\textasciicircum{}T & mode2p\_target\_arrays; mode2n\_source\_target & validated\\
EQ05 & Stokes parameters & S\_0=|E\_x|\textasciicircum{}2+|E\_y|\textasciicircum{}2,\textbackslash{} S\_1=|E\_x|\textasciicircum{}2-|E\_y|\textasciicircum{}2,\textbackslash{} S\_2=2\textbackslash{}Re(E\_xE\_y\textasciicircum{}*),\textbackslash{} S\_3=-2\textbackslash{}Im(E\_xE\_y\textasciicircum{}*) & VectorField.stokes; stokes\_from\_linear\_components & validated\\
EQ06 & HWP synthesis & J\_\textbackslash{}mathrm\{HWP\}(\textbackslash{}beta)[1,0]\textasciicircum{}T=[\textbackslash{}cos(2\textbackslash{}beta),\textbackslash{}sin(2\textbackslash{}beta)]\textasciicircum{}T,\textbackslash{}quad \textbackslash{}beta=\textbackslash{}alpha/2 & route\_patterned\_hwp\_ideal & validated in simulation\\
EQ07 & QWP matrix & J(\textbackslash{}delta,\textbackslash{}psi)=R(-\textbackslash{}psi)\textbackslash{}operatorname\{diag\}(e\textasciicircum{}\{-i\textbackslash{}delta/2\},e\textasciicircum{}\{i\textbackslash{}delta/2\})R(\textbackslash{}psi) & linear\_retarder; apply\_uniform\_jones & model-supported; mount sign calibration-required\\
EQ08 & Dual-SLM phase convention & \textbackslash{}phi\_H=+\textbackslash{}alpha,\textbackslash{}quad \textbackslash{}phi\_V=-\textbackslash{}alpha+\textbackslash{}pi/2 & route\_dual\_slm\_linear\_then\_qwp\_ideal; run\_mode2q\_backward\_initialisation & validated\\
EQ09 & Sequential SLM chain & E\_4=A/\textbackslash{}sqrt2[e\textasciicircum{}\{i\textbackslash{}phi\_H\},e\textasciicircum{}\{i\textbackslash{}phi\_V\}]\textasciicircum{}T & sequential\_jones\_equivalence & validated in simulation\\
EQ10 & Axicon transverse wavevector & k\_r=(2\textbackslash{}pi/\textbackslash{}lambda)(n\_\textbackslash{}mathrm\{axicon\}-n\_m)\textbackslash{}tan\textbackslash{}gamma & NathanSourceParityConfig.k\_r\_m\_inv & validated against source branch\\
EQ11 & Angular spectrum propagation & \textbackslash{}tilde\{\textbackslash{}mathbf E\}(k\_x,k\_y,z)=\textbackslash{}tilde\{\textbackslash{}mathbf E\}(k\_x,k\_y,0)e\textasciicircum{}\{ik\_z z\},\textbackslash{}quad k\_z=\textbackslash{}sqrt\{k\textasciicircum{}2-k\_x\textasciicircum{}2-k\_y\textasciicircum{}2\} & propagate\_vector\_asm & validated in simulation\\
EQ12 & Inverse propagation & \textbackslash{}tilde\{\textbackslash{}mathbf E\}(k\_x,k\_y,0)=\textbackslash{}tilde\{\textbackslash{}mathbf E\}(k\_x,k\_y,z)e\textasciicircum{}\{-ik\_z z\} & mode2q\_backpropagate\_vector & validated in simulation\\
EQ13 & 4F carrier hologram & \textbackslash{}phi\_\textbackslash{}mathrm\{display\}=\textbackslash{}phi\_\textbackslash{}mathrm\{target\}+2\textbackslash{}pi\textbackslash{}nu x & run\_mode2n\_dual\_slm\_4f\_route; mode2s\_apply\_4f & validated in simulation\\
EQ14 & Fourier-order displacement & x\_\{+1\}=\textbackslash{}lambda f\textbackslash{}nu & physical\_4f\_rows; fourf\_final\_design & model-supported; bench calibration-required\\
EQ15 & Power observable & I=|E\_x|\textasciicircum{}2+|E\_y|\textasciicircum{}2+|E\_z|\textasciicircum{}2 & VectorField.intensity & validated simulation observable\\
EQ16 & Ring-count relation & N\_\textbackslash{}mathrm\{rings\}=R k\_r/(2\textbackslash{}pi) & effective\_ring\_count\_for\_plane & audited after plane-mixing correction\\
EQ17 & Strict symmetry diagnostics & \textbackslash{}\{c\_\{60\},c\_\{90\},c\_\{120\},h\_3,h\_4,h\_6,D\_\textbackslash{}mathrm\{core\}\textbackslash{}\} & mode2q\_strict\_hexagon\_gate; evaluate\_strict\_hexagon\_metrics & validated against truth table\\
\bottomrule
\end{longtable}

\chapter{Claim Registry}

\begin{longtable}{p{0.18\textwidth}p{0.18\textwidth}p{0.18\textwidth}p{0.18\textwidth}p{0.18\textwidth}}
\caption{Claim and evidence registry excerpt.}\label{tab:claim-registry}\\
\toprule
claim\_id & claim & numerical\_result & status & limitations\\
\midrule
\endfirsthead
\toprule
claim\_id & claim & numerical\_result & status & limitations\\
\midrule
\endhead
CL01 & V0 literal source and project implementation reproduce the same source-scale field. & source parity validated; Nathan axis-sampled grid required & validated & Simulation parity, not an experimental measurement.\\
CL02 & The ideal patterned-HWP and dual-SLM/QWP Jones routes reproduce the target vector field. & component overlaps equal to 1 within numerical precision & validated & Ideal component model before carrier, 4F and axicon propagation.\\
CL03 & The realistic source-scale dual-SLM + carrier + common 4F route preserves the strict hexagon. & realistic sequential z60 correlation approx 0.9936209493; strict class visual\_hexagonal\_field & validated in simulation & Experimental replication remains unverified.\\
CL04 & M2Q inverse propagation independently recovers the same H/V mask convention. & complex vector overlap approx 0.99999997; alpha RMS approx 1.1e-4 rad & validated in simulation & Inverse uses the simulated complex V0 target.\\
CL05 & The source-scale system is sensitive primarily to hologram/mask centre versus axicon-axis registration. & 0.5 mm decentre recovered from corr approx 0.7888 to approx 0.9762 & model-supported & Finite perturbation set; bench correction still required.\\
CL06 & N=1536 publication-selected outputs are used for the primary hero comparison where available. & Fig. 3A native N=1536; SAS zoom dx approx 3.087 um & validated & Some broad propagation context remains N=1024 by design.\\
CL07 & The old M2U2 best-compromise optimiser candidate is forbidden. & m2u2\_opt\_003\_c5.75\_i0.32\_q-0.25\_r+0.0\_p0.10 & forbidden & Forbidden within this project branch and gate definition.\\
CL08 & The original MODE 2W split-arm presentation is superseded and must not be used as canonical architecture. & M2W superseded; M2WF-A accepted & superseded & Historical split-arm files retained only as provenance.\\
CL09 & Exact SLM phase stroke and per-panel LUT remain calibration-required. & phase stroke unresolved at 1029 nm & calibration-required & Bench calibration required before hardware masks are final.\\
CL10 & Microfabrication/sample-plane success is not claimed by this report. & no microfabrication/sample-plane success claim & validated scope boundary & Separate branch remains unresolved.\\
\bottomrule
\end{longtable}

\chapter{Number Registry}

\begin{longtable}{p{0.18\textwidth}p{0.18\textwidth}p{0.18\textwidth}p{0.18\textwidth}p{0.18\textwidth}}
\caption{Master numerical registry excerpt.}\label{tab:number-registry}\\
\toprule
number\_id & quantity & value & units & status\\
\midrule
\endfirsthead
\toprule
number\_id & quantity & value & units & status\\
\midrule
\endhead
N001 & Physical bench wavelength & 1029 & nm & resolved\\
N002 & Nathan parity wavelength & 1030 & nm & different scope\\
N003 & Source beam radius & 2.0 & mm & validated\\
N004 & SLM panel width & 1920 & px & repository-confirmed geometry\\
N005 & SLM panel height & 1080 & px & repository-confirmed geometry\\
N006 & SLM pitch & 8 & um & repository-confirmed geometry\\
N007 & SLM active area & 15.36 x 8.64 & mm & repository-confirmed geometry\\
N008 & Carrier frequency & 6.25 & lp/mm & validated\\
N009 & Carrier period & 20 & SLM px & validated\\
N010 & 4F focal length & 300 & mm & routine bench confirmation\\
N011 & +1 order displacement & 1.929 & mm & model-supported\\
N012 & Iris diameter & 1.54 & mm & routine bench calibration\\
N013 & Axicon base angle & 2 & deg & validated\\
N014 & Axicon refractive index & 1.458 & dimensionless & validated\\
N015 & Reference z & 60 & mm & validated\\
N016 & Publication grid & 1536 & samples per side & publication-selected\\
N017 & N=1536 samples per radial fringe & 9.89 & samples/fringe & validated\\
N018 & Fig. 3A SAS zoom sampling & 20.84 & samples/fringe & display provenance\\
N019 & Realistic sequential z60 correlation & 0.9936209493 & correlation & validated in simulation\\
N020 & Sequential pre-QWP overlap & 1.0 & overlap & validated in simulation\\
N021 & Sequential post-QWP overlap & 1.0 & overlap & validated in simulation\\
N022 & First-order efficiency & 0.949 & fraction & validated in simulation\\
N023 & 0.5 mm compensated axicon decentre correlation & 0.9762 & correlation & model-supported\\
N024 & Old forbidden candidate full-field correlation & 0.989 & correlation & forbidden despite high correlation\\
\bottomrule
\end{longtable}

\chapter{Superseded and Forbidden Material}

\begin{longtable}{p{0.18\textwidth}p{0.18\textwidth}p{0.18\textwidth}p{0.18\textwidth}p{0.18\textwidth}}
\caption{Superseded and invalidated material.}\label{tab:superseded}\\
\toprule
item\_id & material & reason\_superseded & replacement & status\\
\midrule
\endfirsthead
\toprule
item\_id & material & reason\_superseded & replacement & status\\
\midrule
\endhead
S01 & Original M2W split-arm optical diagram & It depicted a PBS split/recombine interferometer as canonical; M2W-FIX accepts a sequential single-beam route. & F08 sequential optical architecture and docs/83 & superseded\\
S02 & Split-arm H/V/PBS power ledger & Sequential route has no separate H/V spatial arms and no PBS recombination power bookkeeping. & F14 and mode2w\_fix\_sequential\_power\_ledger.csv & superseded\\
S03 & m2u2\_opt\_003\_c5.75\_i0.32\_q-0.25\_r+0.0\_p0.10 & High full-field correlation hid fourfold/X-like non-hexagonal structure. & REALISTIC\_4F\_HEXAGON\_REFERENCE & forbidden\\
S04 & Low-resolution hero montages & High DPI/interpolation cannot replace numerical sampling. & N=1536 / SAS-audited M2W-FIX hero figures & superseded\\
\bottomrule
\end{longtable}

\chapter{Figure Manifest}

\begin{longtable}{p{0.18\textwidth}p{0.18\textwidth}p{0.18\textwidth}p{0.18\textwidth}p{0.18\textwidth}}
\caption{Figure provenance manifest excerpt.}\label{tab:fig-manifest}\\
\toprule
figure\_id & report\_caption & source\_stage & N & validation\_status\\
\midrule
\endfirsthead
\toprule
figure\_id & report\_caption & source\_stage & N & validation\_status\\
\midrule
\endhead
F01 & V0 source parity and reference/reproduction comparison. & V0 & 1024 & validated\\
F02 & Grid convention audit showing why Nathan axis-sampled centring is required. & V0 & 1024 & validated\\
F03 & High-resolution V0 reference crop and propagated structure. & M2U & 1536 & validated\\
F04 & M2P target alpha and sector map. & M2P/M2U & 1536 & validated\\
F05 & M2P pre-axicon Stokes comparison. & M2P/M2U & 1536 & validated\\
F06 & M2Q recovered phase-only masks and inverse design diagnostics. & M2Q/M2U & 1536 & validated\\
F07 & M2W-FIX V0 versus ideal sequential versus realistic sequential comparison. & M2W-FIX & 1536 & validated\\
F08 & Corrected sequential single-beam optical system. & M2W-FIX & diagram & validated architecture\\
F09 & Native SLM1/SLM2 masks and target vector-field diagnostics. & M2W-FIX & 1920x1080 masks; 1536 target grid & validated display package\\
F10 & Realism degradation: clean, moderate and bad cases. & M2W-FIX/M2S & 1024 & model-supported\\
F11 & Corrected recovery from a 0.5 mm axicon/mask offset. & M2W-FIX/M2S/M2V & 1024 & model-supported\\
F12 & Transverse propagation with physical context and SAS-scaled focus crops. & M2W-FIX & 1024 & validated display provenance\\
F13 & x-z and y-z propagation maps plus z diagnostics. & M2W-FIX & 1024 & validated\\
F14 & Sequential-route power flow. & M2W-FIX & ledger & model-supported\\
F15 & Tolerance limits table/plot. & M2W-FIX/M2S & mixed & model-supported\\
F16 & Combined realism and correction results. & M2W-FIX/M2S/M2V & mixed & model-supported\\
F17 & Strict classifier truth table and gate calibration. & M2U2-FIX & highres & validated gate\\
S01 & Forbidden old best-compromise audit figure. & M2U2-FIX & historical & forbidden/superseded\\
\bottomrule
\end{longtable}

\end{document}